% Lösungen Einheit 1 von 4 – Gleichungen verstehen (zusammenbau v0.1)
\einheitenkopf{Einheit 1 von 4 · Gleichungen verstehen}

\erg{8}{a) $-6$}
\erg{9}{a) $6 + 9 = 15$ (wA) \quad b) $2 \cdot 3 + 7 = 13$ (wA) \quad c) $3 \cdot 8 + 4 = 28$, nicht 25 (fA) \quad d) links $5 \cdot 4 - 2 = 18$, rechts $3 \cdot 4 + 6 = 18$ (wA) \quad e) $3x + 5 = -4$, denn $3 \cdot (-3) + 5 = -4$ (wA)}
\erg{10}{a) $x = 4$ (Tabelle: 5, 7, 9, 11) \quad b) $x = 10$ \quad c) $x = 8$, denn $21 - 13 = 8$ \quad d) $4$, denn $7 \cdot 4 - 6 = 22$ ist eine wahre Aussage \quad e) $x = 36$; Probe: $36 + 47 = 83$ (wA)}
\erg{11}{a) Er hat die rechte Seite nicht ausgerechnet: rechts $4 + 10 = 14$, links 14 – 4 ist Lösung (wA). \quad b) Die Probe fehlt: $4 \cdot 9 + 6 = 42$, nicht 30 (fA). Richtig ist $x = 6$: $4 \cdot 6 + 6 = 30$ (wA). \quad c) Nein: Eine positive Zahl plus 25 ist größer als 25, also nie 10. \quad d) $x - 6 = 15$ \quad e) „Welche Zahl plus 14 ergibt 30?" – $x = 16$}
